% experiments.tex -- Module 7 ESSENTIALS: Key experiments at a glance.
% One row per Key experiment box of ch1.tex ... ch4.tex in this folder, plus the historical
% experiments that those chapters give as short bullets (Hertz, Romer, the microwave oven, the
% 1972 laser measurement, Foucault, Young, the muon experiments, Bertozzi). Facts and numbers as in those chapters, which follow the full book
% (src/mod7/ch*.tex). "Ch" links to the section that describes the experiment.
\unnumberedchapter[\faFlask]{Key experiments at a glance}

The experiments of this module, in chapter order (``Ch'': click to go there).
\textbf{(inv.)} = an investigation: practical work you may have done in class.

\par\smallskip
{\small
\setlength\LTpre{4pt}\setlength\LTpost{4pt}
\renewcommand\arraystretch{1.12}\arrayrulecolor{hscrule}
\newcommand\essexhead{\colhead{Experiment} & \colhead{What was done} &
  \colhead{What it showed} & \colhead{Ch}}
% \essexch{section label}{chapter number}
\newcommand\essexch[2]{\hyperref[#1]{\textcolor{hscbrand}{\bfseries #2}}}
\begin{longtable}{@{}L{30mm}@{\hspace{3mm}}L{52mm}@{\hspace{3mm}}%
    L{\dimexpr\linewidth-30mm-52mm-5mm-9mm\relax}@{\hspace{3mm}}C{5mm}@{}}
  \essexhead\\\hline\endfirsthead
  \essexhead\\\hline\endhead
  \textbf{Hertz} (1887)
    & made and detected radio waves with sparks
    & they acted like light, at about $c$: Maxwell was right
    & \essexch{ch1:sec:maxwell}{1}\\\hline
  \textbf{Rømer} (1676)
    & Io's eclipses timed as Earth moved away from Jupiter
    & up to \qty{22}{min} late: light is not instant
    & \essexch{ch1:sec:speed}{1}\\\hline
  \textbf{Fizeau's wheel} (1849)
    & a spinning 720-tooth wheel; a mirror \qty{8.63}{km} away
    & $c \approx \qty{3.1e8}{m.s^{-1}}$: the first measurement of $c$ on Earth
    & \essexch{ch1:sec:experiments}{1}\\\hline
  \textbf{Microwave oven} (inv.)
    & melted chocolate spots in a standing wave; $c = f\lambda$
    & $c \approx \qty{3.0e8}{m.s^{-1}}$ (an indirect method)
    & \essexch{ch1:sec:speed}{1}\\\hline
  \textbf{Laser $f$ and $\lambda$} (1972)
    & $f$ and $\lambda$ of a laser measured; $c = f\lambda$
    & $c$ to 9 figures; since 1983 $c$ is fixed and defines the metre
    & \essexch{ch1:sec:speed}{1}\\\hline
  \textbf{Spectra} (inv.)
    & gas tubes, a filament lamp and sunlight through a spectroscope
    & each gas has its own lines; lamp: continuous; Sun: dark lines
    & \essexch{ch1:sec:experiments}{1}\\\hline
  \textbf{Foucault} (1850)
    & spinning-mirror timing, with and without water in the path
    & light is \emph{slower} in water: Newton was wrong
    & \essexch{ch2:sec:models}{2}\\\hline
  \textbf{Young} (about 1801)
    & light from one pinhole through two narrow slits close together
    & bright and dark fringes: interference, so light is a wave
    & \essexch{ch2:sec:interference}{2}\\\hline
  \textbf{Lasers} (inv.)
    & lasers through single slits, a double slit and a grating
    & narrower slit or longer $\lambda$: more spreading; $\lambda$ within about 0.3\,\%
    & \essexch{ch2:sec:experiments}{2}\\\hline
  \textbf{Malus' law} (inv.)
    & an analyser turned in \ang{15} steps after a polariser; a light sensor
    & $I \propto \cos^2\theta$: light is a transverse wave
    & \essexch{ch2:sec:experiments}{2}\\\hline
  \textbf{Black body curves} (1899--1900)
    & oven, prism and bolometer, at several temperatures
    & Wien's law; classical theory failed, so Planck used quanta
    & \essexch{ch3:sec:experiments}{3}\\\hline
  \textbf{Photoelectric effect} (1902--1916)
    & $V_0$ for several frequencies and brightnesses (Lenard, Millikan)
    & $V_0$ rises in a straight line with $f$: photons; $h$ from the gradient
    & \essexch{ch3:sec:experiments}{3}\\\hline
  \textbf{Michelson and Morley} (1887)
    & an interferometer, floating on mercury, slowly rotated
    & a shift under 0.02 of a fringe, not 0.4: no aether wind
    & \essexch{ch4:sec:experiments}{4}\\\hline
  \textbf{Muons} (1963, 1977)
    & counted on a mountain and at sea level; stored in a ring
    & far more survive; ring muons live 29.3 times longer: time dilation
    & \essexch{ch4:sec:time}{4}\\\hline
  \textbf{Hafele and Keating} (1971)
    & caesium atomic clocks flown east and west around the world
    & time differences as predicted: moving clocks run slow
    & \essexch{ch4:sec:experiments}{4}\\\hline
  \textbf{Bertozzi} (1964)
    & electrons given up to \qty{15}{\mega\eV}, then timed
    & they came close to $c$ but never passed it
    & \essexch{ch4:sec:momentum}{4}\\\hline
  \textbf{Supernovae}
    & distant Type Ia supernovae timed as they fade
    & they fade more slowly by the factor $1 + z$: time is not absolute
    & \essexch{ch4:sec:experiments}{4}\\\hline
\end{longtable}}
